\DocumentMetadata{
  pdfversion=2.0,
  pdfstandard=ua-2,
  lang=en-US,
  tagging=on,
  tagging-setup={math/setup=mathml-SE,tabsorder=structure}
}
\documentclass[11pt,letterpaper]{article}
\usepackage[margin=0.68in,includefoot]{geometry}
\usepackage{newcomputermodern}
\usepackage{amsmath,amscd,mathtools}
\providecommand{\square}{\mdlgwhtsquare}
\usepackage[hidelinks]{hyperref}
\hypersetup{
  pdftitle={Numerical Analysis PhD Exam, August 1995},
  pdfauthor={University of Florida Department of Mathematics},
  pdfsubject={Graduate mathematics examination},
  pdfdisplaydoctitle=true,
  bookmarksopen=true
}
\setcounter{secnumdepth}{0}
\setlength{\parindent}{0pt}
\setlength{\parskip}{0.28em}
\setlength{\emergencystretch}{3em}
\setlength{\leftmargini}{2.15em}
\setlength{\leftmarginii}{2.65em}
\setlength{\leftmarginiii}{3em}
\renewcommand{\labelenumi}{\textbf{\theenumi.}}
\pagestyle{empty}
\raggedbottom
\allowdisplaybreaks
\newenvironment{examproblems}[1]{%
  \begin{enumerate}%
  \setcounter{enumi}{#1}%
  \setlength{\topsep}{0.35em}%
  \setlength{\partopsep}{0pt}%
  \setlength{\itemsep}{0.48em}%
  \setlength{\parsep}{0.18em}%
}{\end{enumerate}}
\newenvironment{examsubparts}{%
  \begin{enumerate}%
  \setlength{\topsep}{0.18em}%
  \setlength{\partopsep}{0pt}%
  \setlength{\itemsep}{0.16em}%
  \setlength{\parsep}{0.1em}%
}{\end{enumerate}}
\newenvironment{examinstructions}{%
  \par\begingroup\itshape
}{%
  \par\endgroup\vspace{0.18em}
}
\makeatletter
\renewcommand\section{\@startsection{section}{1}{0pt}%
  {-1.25ex plus -.25ex minus -.1ex}%
  {0.55ex plus .1ex}%
  {\normalfont\large\bfseries}}
\renewcommand{\@maketitle}{%
  \newpage\null\vskip -1.2em%
  {\footnotesize\bfseries
    \href{https://www.ufl.edu/}{University of Florida}, \href{https://math.ufl.edu/}{Department of Mathematics}\par}%
  \vskip 0.18em%
  \tagstructbegin{tag=Title}%
    {\LARGE\bfseries\href{https://gma.math.ufl.edu/exams/numerical-analysis-phd/}{Numerical Analysis PhD Exam}, August 1995\par}%
  \tagstructend%
  \vskip 0.35em%
  \tagmcbegin{artifact}%
    \hrule height 1pt%
  \tagmcend%
  \vskip 0.55em%
}
\makeatother
\title{Numerical Analysis PhD Exam, August 1995}
\author{}
\date{}
\begin{document}
\maketitle
\thispagestyle{empty}
\begin{examproblems}{0}
\item
Consider the Householder matrix \(H=I-2ww^T\), where \(w^Tw=1\).
\begin{examsubparts}
\item[\textnormal{(a)}]
Show that~\(H\) is orthogonal.
\item[\textnormal{(b)}]
Determine the eigenvalues of~\(H\).
\item[\textnormal{(c)}]
Given a vector~\(x\) and an integer~\(k\), find a vector~\(w\) such that \((Hx)_i=x_i\) for \(i<k\) and \((Hx)_i=0\) for \(i>k\).
\item[\textnormal{(d)}]
Explain why one needs to be careful about the choice of sign in the formula for~\(w\).
\end{examsubparts}
\item
Let~\(A\) be a symmetric, positive definite matrix. Show that if~\(A\) is \(LU\) factored, then the elements of~\(U\) in absolute value are bounded by the largest diagonal element of~\(A\).
\item
Give a careful statement and proof of Gershgorin’s Theorem. Find upper and lower bounds for the set of eigenvalues for the matrix 
\[
\begin{pmatrix}
2&-1&0&0\\
-1&2&-1&0\\
0&-1&2&-1\\
0&0&-1&2
\end{pmatrix}.
\]
\item
Let~\(A\) be a nonsingular \(n\times n\) matrix and \(\{X_k\}\) a sequence of \(n\times n\) matrices generated by the iteration 
\[
X_{k+1}=X_k+X_k(I-AX_k).
\]
 Show that if \(X_0\) is chosen such that the spectral radius of \(I-AX_0\) is strictly less than 1, then \(\{X_k\}\) converges to the inverse of~\(A\).
\item
For a given partition \(\{x_0,x_1,\ldots,x_n\}\) of an interval \([a,b]\) and function~\(f\) defined on \([a,b]\):
\begin{examsubparts}
\item[\textnormal{(a)}]
State the divided difference formula for the polynomial \(p_n(x)\) interpolating~\(f\) from the data \((x_k,f(x_k))\), \(k=0,1,\ldots,n\).
\item[\textnormal{(b)}]
Show that the error in the approximation in part (a) is given by 
\[
f(x)-p_n(x)=(x-x_0)(x-x_1)\cdots(x-x_n)f[x_0,x_1,\ldots,x_n,x].
\]
\item[\textnormal{(c)}]
For \(\{x_0,x_1,x_2\}=\{0,1,2\}\) and \(f(x)=1/(1+x^2)\), compute the divided difference formula for \(p_2(x)\).
\end{examsubparts}
\item
Let~\(A\) be a real, positive definite \(n\times n\) matrix and~\(b\) be an \(n\times1\) vector.
\begin{examsubparts}
\item[\textnormal{(a)}]
Derive the steepest descent algorithm for approximating the true solution~\(x\) of \(Ax=b\), by a sequence of iterates \(\{x_k\}\).
\item[\textnormal{(b)}]
Define the real-valued function~\(h\) on \(\mathbb{R}^n\) by 
\[
h(x)=(Ax-b)^TA^{-1}(Ax-b)
\]
 and show that 
\[
h(x_{k+1})\leq\left(1-c(A)^{-1}\right)^2h(x_k),
\]
 where \(c(A)\) is the condition number of the matrix~\(A\).
\end{examsubparts}
\item
For an algorithm, define the terms inherent error, total effect of rounding, and numerically stable.

\par
Carefully formulate a numerically stable algorithm for computing the positive solution of the quadratic equation \(x^2+2px-q=0\), where~\(p\) and~\(q\) are arbitrary positive numbers. Justify the numerical stability of your algorithm.
\item
Give a careful proof that the DFT of a convolution of two finite sequences (not necessarily of the same length) is proportional to the product of the DFT’s of the original sequences, suitably padded.
\item
Given the~\(N\) sampled values \(f(2k\pi/N)\), \(k=0,1,2,\ldots,N-1\), of a twice continuously differentiable, \(2\pi\)-periodic function~\(f\), outline a method of computing the Fourier coefficients of~\(f\) which has the following properties:

\par
• makes use of the FFT routine;

\par
• produces an estimate which improves as \(N\to\infty\).

\par
Estimate the accuracy of your method. If possible, state an explicit formula for its implementation.
\end{examproblems}
\end{document}
